% Ce qui reste à faire, écrit à la main à partir de v2/docs/lab/open.md (99 questions ouvertes au
% 14 novembre 2026) et des rapports d'axe. Coûts en temps de travail estimé (\suppose), sauf mention.
\part{Ce qui reste à faire}

\section{Par axe}
\paragraph{Axe 0, le jumeau.}
\begin{itemize}
  \item \textbf{Vitesse à la commune.} Il faut environ 6,8~s par scénario (extrapolé) pour une
        cible de 2~s, ce qui bloque le zoom de l'ARS jusqu'à la commune. Coût : environ 2~h.
  \item \textbf{Vivier réel.} Utiliser les effectifs réels en formation par profession, au
        lieu d'un vivier relatif supposé.
  \item \textbf{Diplômés étrangers} : un âge d'entrée propre.
  \item \textbf{Obligation d'exercice} : à brancher quand le modèle d'installation de l'axe~2
        existera.
\end{itemize}

\paragraph{Axe 1, les hyperparamètres.}
\begin{itemize}
  \item \textbf{Lire en texte intégral toutes les sources citées}, et vérifier chaque nombre.
        L'état de l'art méthodologique n'a été lu qu'en résumé.
  \item \textbf{Trouver un flux annuel de diplômés étrangers.} On ne connaît aujourd'hui que
        des parts de stock : 24~\% des dentistes et 33~\% des kinés n'ont pas de diplôme
        d'État français (\mesure).
  \item \textbf{Expliquer l'écart sur la retraite} : le jumeau donne $+3$~\% chez les
        infirmiers, contre environ $+1$~\% cité par la DREES.
  \item \textbf{Re-prioriser le sourçage} à partir des indices de l'axe~5.
\end{itemize}

\paragraph{Axe 2, les méthodes.}
\begin{itemize}
  \item \textbf{Projection communale : arrêtée en attente des données 2027.} Les modèles communaux K3, K4 et K5,
        chacun inscrit avant son passage, ont échoué au critère N10 (couverture de 75 à 85~\% dans chaque classe de
        taille) ou à leur confirmation hors échantillon. L'affichage reste au département ; les dépôts K4, K5 et C1
        seront notés sur les effectifs 2027.
  \item \textbf{Étape 3} : un logit conditionnel d'installation, qui attend les flux
        communaux de l'axe~4.
  \item \textbf{Étapes 4 et 5} : la réconciliation MinT entre échelles, et la calibration par
        simulation sur les cinq professions.
\end{itemize}

\paragraph{Axe 3, la validation.}
\begin{itemize}
  \item \textbf{Exécuter N6} : CRPS, PIT et couverture, contre une référence naïve elle aussi
        probabiliste.
  \item \textbf{Calibration par simulation} sur 200 mondes simulés.
  \item \textbf{N1 sur plusieurs origines}, pour vérifier la stabilité du classement des
        estimateurs.
  \item \textbf{Dentistes : écart signalé, puis levé.} Le chiffre de 76\,000 est la projection
        tendancielle \emph{2050}, pas l'effectif 2026 : en 2026, le jumeau compte 47\,347 dentistes
        dans le privé, sous le total RPPS tous modes de 48\,690 au 1\textsuperscript{er} janvier 2026
        (\mesure ; rapport \texttt{2026-11-14-reconciliation-dentistes.md}). Le chiffre de 62\,000
        attribué à la DREES n'a pas de source retrouvée.
\end{itemize}

\paragraph{Axe 4, les données.}
\begin{itemize}
  \item \textbf{Installations par commune.} Archiver chaque mois l'Annuaire Santé, la seule voie
        ouverte vers des flux individuels. En attendant, construire un substitut à partir de
        CartoSanté : les installations nettes des moins de 40~ans, de 2004 à 2025.
  \item \textbf{Officines} : ouvertures et fermetures par commune, d'après l'historique FINESS.
  \item \textbf{Trous restants} : pas d'historique communal pour les sages-femmes, et pas de
        croisement âge $\times$ secteur pour les infirmiers.
\end{itemize}

\paragraph{Axe 5, la sensibilité.}
\begin{itemize}
  \item \textbf{Amplitude de référence fixée} à celle de la DREES (quotas $\pm 20$~\%) : en 2050, la décision pèse
        plus que l'incertitude pour les infirmiers (0,89), les kinés (0,81) et les sages-femmes (0,59), à égalité pour
        les dentistes, moins pour les pharmaciens (0,38 contre 0,61).
  \item \textbf{Étendre les indices de Sobol aux 12 leviers branchés}, pour séparer la part de
        variance qui vient de la décision de celle qui vient de notre ignorance.
  \item \textbf{Cartographier les fourchettes 2050 par département.}
  \item \textbf{Classer les sources d'incertitude} selon le coût pour les réduire.
\end{itemize}

\section{Pour l'interface}
Faits depuis les tests : repères 2030 à 2050 dans la fiche, fiche exportable avec sources et hypothèses, densité pour
100\,000 habitants, phrase en langage naturel, mode expert. Restent :
\begin{enumerate}
  \item une \textbf{recette rejouée sur la carte finale}, un persona par onglet, captures sur disque ;
  \item la \textbf{densité au grain infra-départemental}, à partir de la population communale ;
  \item l'\textbf{écran d'accueil} « Quel territoire ? », avec les modes simple et expert ;
  \item des \textbf{libellés sans jargon}, et des erreurs réseau signalées ;
  \item les \textbf{tokens W3C}, pour reprendre la direction artistique dans Figma.
\end{enumerate}

\section{Décisions qui reviennent à l'équipe}
\begin{itemize}
  \item \textbf{Compilateur : modèle de langage ou non ?} L'approche à glossaire plafonne vers 77~\% ; la décision
        prévoit, dans ce cas, un modèle de langage pour l'analyse de la phrase seule (jamais le calcul), évalué sur
        les mêmes jeux écrits à l'aveugle.
  \item \textbf{Un module de besoin ?} Les leviers de transfert de tâches (vaccination en
        officine, infirmiers en pratique avancée, accès direct au kiné) ne changent pas
        l'effectif. Ils changent le besoin auquel il répond. C'est pourtant l'exemple même du
        sujet : « combien de spécialistes en 2035 si les pharmaciens peuvent vacciner ? ».
        Sans module de besoin, ces leviers n'ont aucun effet dans le jumeau.
  \item \textbf{Archiver l'Annuaire Santé chaque mois} : c'est une tâche planifiée, qui
        demande l'accord de l'équipe.
  \item \textbf{Garder les quatre familles de leviers} (former, attirer, installer, retenir),
        ou en choisir d'autres.
  \item \textbf{Phase 2} (cardiologie, dermatologie, ophtalmologie) : à démarrer une fois les
        axes 0 à 3 terminés pour au moins une profession de la phase 1.
\end{itemize}
